\begin{figure*}[t]
\centering
\scriptsize
\setlength{\fboxsep}{1em}
\setlength{\fboxrule}{1pt}
\fcolorbox{black}{white}{
\parbox{0.955\textwidth}{
\begin{alltt}


You will be provided instructions for completing a task followed by a task to complete.

Instructions:

1. Understand the question: Read the question carefully and understand what it's asking. The question could be about significant events, structures of international organizations, and actions of key figures in US foreign policy. 

2. Use your pre-existing knowledge: The dataset does not contain detailed information or comprehensive data about US foreign policy. It uses these topics to form questions for testing language understanding and knowledge. Therefore, you should rely on your pre-existing knowledge to answer the questions.

3. Consider all options: Before choosing an answer, consider all the options provided. Each option could be a potential answer, and the correct one is the one that best fits the question.

4. Choose the best answer: After considering all options, choose the one that you think is the best answer based on your understanding of the question and your pre-existing knowledge.

5. Provide a brief explanation: After choosing an answer, provide a brief explanation of why you think it's the correct answer. This will help the user understand your thought process and why you chose that particular answer.

6. Check your answer: Before finalizing your answer, check it again to make sure it's the best answer to the question. If you're not sure, consider the question and options again, and see if there's a better answer.

Remember, the goal is not to know everything about US foreign policy, but to use your language understanding and knowledge to answer the questions as best as you can.

Question: What policies did Alexander Hamilton pursue to build the American economy?

A. Free trade with all European countries on an equal basis and a smaller government role in the economy

B. Establishment of a national currency and higher government spending

C. Tariffs to protect American manufacturing, establishment of a national bank, and the repayment of all debts.
D. Limits on immigration and imported technology to encourage higher wages

Follow the instructions to generate an explanation that reasons towards the correct answer to the task above. End the explanation with the correct answer. [/INST] Explanation: 

To answer this question, we need to consider the policies that Alexander Hamilton pursued during his tenure as the first Secretary of the Treasury. Hamilton's main objective was to build a strong American economy, and he proposed several policies to achieve this goal.

Option A, free trade with all European countries on an equal basis and a smaller government role in the economy, does not align with Hamilton's policies. Hamilton believed in a strong central government and advocated for protectionist policies to promote American industry.

Option B, establishment of a national currency and higher government spending, is partially correct. Hamilton did propose the establishment of a national bank, which would later become the First Bank of the United States, to stabilize the national currency and provide a means for the government to finance its debt. However, Hamilton did not advocate for higher government spending.

Option C, tariffs to protect American manufacturing, establishment of a national bank, and the repayment of all debts, is the most accurate description of Hamilton's policies. Hamilton believed that a strong manufacturing sector was essential for the country's economic growth and proposed tariffs to protect American industries from foreign competition. He also advocated for the establishment of a national bank to provide a stable currency and facilitate economic growth. Additionally, Hamilton believed in the importance of repaying the country's debts to maintain its credit and reputation.

Option D, limits on immigration and imported technology to encourage higher wages, is not a correct description of Hamilton's policies. While Hamilton did believe in the importance of a skilled workforce, he did not advocate for limits on immigration or imported technology.

Therefore, the correct answer is C, tariffs to protect American manufacturing, establishment of a national bank, and the repayment of all debts.

Therefore, the correct multiple choice label (just the letter) to the task is below.

Answer: C
\end{alltt}
}}
\caption{A complete running example of MMLU (US Foreign Policy).}
\label{fig:mmlu_subject=us_foreign_policy_example}
\end{figure*}
